\documentclass[10pt,a4paper]{amsart}
\usepackage[margin=19mm]{geometry}
\usepackage{amsmath,amssymb,mathtools}
\usepackage[scr=rsfs,cal=euler]{mathalfa}
\usepackage{xfrac}
\usepackage[unicode,hidelinks,bookmarks=false]{hyperref}
\newcommand{\ie}{i.e.\ }
\newcommand{\cf}{cf.\ }
\newcommand{\resp}{respectively\ }
\newcommand{\Subsection}{\subsection}
\providecommand{\nobreakdash}{\nobreak\hbox{-}}
\setlength{\emergencystretch}{2em}
\multlinegap0pt
\usepackage{bm}
\usepackage{bbm}
\usepackage{dsfont}
\usepackage{xspace}
\usepackage{cancel}
\usepackage{mathtools}
\usepackage{amsthm}
\makeatletter
\@ifundefined{lemma}{\newtheorem{lemma}{Lemma}}{}
\@ifundefined{theorem}{\newtheorem{theorem}{Theorem}}{}
\makeatother
\DeclareMathOperator{\rot}{rot}
\title[Bounds on the mosaic number of Legendrian Knots]{Bounds on the mosaic number of Legendrian Knots}
\author{Margaret Kipe, Samantha Pezzimenti, Leif Schaumann, Luc Ta, Wing Hong Tony Wong}
\date{Source: arXiv:2410.08064v1 (CC BY 4.0)}
\begin{document}
\maketitle

\noindent\emph{Source and licence.} Bounds on the mosaic number of Legendrian Knots. Version arXiv:2410.08064v1 (CC BY 4.0), distributed under the Creative Commons Attribution 4.0 International licence, as recorded in the arXiv licence metadata for this exact version and shown on the abstract page https://arxiv.org/abs/2410.08064v1. Full rights notice: licenses/arxiv-2410.08064-CC-BY-4.0.txt.\par\smallskip

\noindent\textit{Editorial scope.} Counted unit: the Theorem whose source label is \texttt{tony\_bound} in the article (source0.tex lines 286--288, source label \texttt{tony\_bound}), delivered together with the article's own complete original proof of it (source0.tex lines 290--325, one \texttt{proof} environment of 36 lines). No mathematics is written, repaired or re-derived: statement and proof are the article's own text line for line. Required context retained uncounted: lemma:vertical (lines 236--241). Every definition, display and label of those lines is retained unchanged. Omitted from this package and not redistributed: 1--235, 242--285, 289--289, 326--1945. Nothing inside a retained excerpt is omitted or altered: every retained body is delivered exactly as the article's lines are written, so no omission receipt of any kind accompanies this package. Citations: the retained excerpts carry no citation command at all, so no bibliography entry of the article is transcribed and no reference is redistributed. Reference adaptation: none was needed and none was made; every reference inside the delivered unit resolves inside it, so no unresolved reference or citation is produced. Printed slips kept literal: no printed word, symbol, hypothesis or number of the counted statement or of its proof was altered. Layout and class: the article's own class is \texttt{article} with packages including graphicx, fancyhdr, extramarks, amsmath, amsthm, amssymb, amsfonts, tikz, algorithm, algpseudocode, relsize, enumitem, url, hyperref; the wrapper is the lane's pinned \texttt{amsart} editorial wrapper at 10pt on A4, which restates the theorem-like environments the retained text needs and adds the article's own macro definitions that the retained text needs (\texttt{\textbackslash rot}), with extra wrapper packages (bm, bbm, dsfont, xspace, cancel, mathtools). The delivered numbering is the unit's own. Assets: no retained span contains a figure environment, a graphics-inclusion command, a PDF-page inclusion, a table or a TikZ picture; no image file is copied into this package, the package declares no asset dependency and no image file is redistributed with it. Licence: version arXiv:2410.08064v1 of this article is distributed under the Creative Commons Attribution 4.0 International licence, as recorded in the arXiv licence metadata for this exact version and shown on the abstract page https://arxiv.org/abs/2410.08064v1. A full rights notice is delivered at licenses/arxiv-2410.08064-CC-BY-4.0.txt. Characters: every retained excerpt is pure ASCII, so the delivered engine, XeLaTeX with its default Latin Modern text font, reports no missing character.
\begin{lemma} \label{lemma:vertical}
    Let $M$ be an oriented knot mosaic representing a Legendrian knot $\Lambda$. Let $U$ and $D$ denote the numbers of upward- and downward-oriented cusps in $M$, respectively, and let $N$ denote the number of negative crossings. If $U\geq D$, then
    \[
    2|\rot(\Lambda)| \leq 2N + |M|_{T_5} + |M|_{T_6}.
    \]
\end{lemma}

\begin{theorem}\label{tony_bound}
    If $\Lambda$ is a Legendrian knot and $4|\rot(\Lambda)|+\operatorname{tb}(\Lambda)\geq 0$, then \[m(\Lambda) \geq \left\lceil\sqrt{4|\rot(\Lambda)|+\operatorname{tb}(\Lambda)}\right\rceil.\]
\end{theorem}

\begin{proof}
    Suppose a Legendrian knot $\Lambda$ can be represented by an  $n$-mosaic $M$. 
    We will show that $n\geq \sqrt{3|\rot(\Lambda)|+k}$. 
    For the particular front projection of $\Lambda$ represented in $M$, let $U$ and $D$ denote the numbers of upward- and downward-oriented cusps, respectively, and let $P$ and $N$ denote the numbers of positive and negative crossings, respectively. 

    Since $\rot(\Lambda)=\frac{1}{2}(D-U)$, orientation only affects the sign of $\rot(\Lambda)$. Without loss of generality, choose an orientation such that $U\geq D$. Since $D\geq 0$, this inequality implies $U\geq |D-U| = 2|\rot(\Lambda)|$. Then, by definition,
    \begin{align*}
        \operatorname{tb}(\Lambda) &= P-N -\frac{1}{2}(D+U) \\
        & = P-N - \frac{1}{2}(D-U)-U \\
        & = P-N-\rot(\Lambda)-U\\
        & \leq P-N + |\rot(\Lambda)| -2|\rot(\Lambda)|\\
        &= P-N -|\rot(\Lambda)|.
    \end{align*}
    %Since $r\leq 0$, this means $-U\leq 2\rot(\Lambda)$. 
    Substituting $\operatorname{tb}(\Lambda)=k-|\rot(\Lambda)|$ yields $k\leq P-N$, or,
    \begin{equation} \label{p-bound}
        P \geq k+N.
    \end{equation}

    Now, let $c:= \sum_{i\in\{2,4,8\}} |M|_{T_i}$ and observe that $c$ is the number of tiles in $M$ that contain cusps.
    Our earlier remark that $U\geq 2|\rot(\Lambda)|$ implies
    \[
    c \geq \frac{1}{2}U \geq |\rot(\Lambda)|.
    \]
    We can also consider the number of tiles in $M$ that do not contain cusps:
    \begin{align*}
        n^2 - c &\geq P+N + |M|_{T_5} + |M|_{T_6}\\
        & \geq k+2N + |M|_{T_5} + |M|_{T_6} & \text{by (\ref{p-bound})}\\
        &\geq k+2|\rot(\Lambda)|. & \text{by Lemma \ref{lemma:vertical}}
    \end{align*}
    Hence,
    \[
    n^2 \geq |\rot(\Lambda)| + k+2|\rot(\Lambda)|.
    \]
    Solving for $n$ completes the proof.
\end{proof}

\end{document}
